%% pre/abstract.tex
\begin{resumo}[Abstract]
 \begin{otherlanguage*}{english}
    \begin{flushleft}
        \setlength{\absparsep}{0pt}
        \SingleSpacing  		\imprimirautorabr~~\textbf{\imprimirtitleabstract}.	\imprimirdata.  \pageref{LastPage} p.
        \imprimirtipotrabalhoabs~-~\imprimirinstituicao, \imprimirlocal, 	\imprimirdata.
    \end{flushleft}
    \OnehalfSpacing
   This work investigates whether news-derived event features can improve the risk-adjusted performance of a Forex trading strategy beyond a price-only baseline. Although the title refers to AI-driven sentiment, the executed study tests GDELT coded-event intensity; article-level sentiment scoring was not executed. The question matters because the Forex market is strongly driven by news flow and macroeconomic information, yet the literature on text-based prediction of financial returns has concentrated on equities, leaving the asset class comparatively underexplored. The executed approach builds, on EUR/USD at minute resolution, a hybrid trading pipeline that combines technical signals with GDELT event-intensity features, fused into a single trade decision under strict point-in-time discipline and a conservative publication-to-use delay. Evaluation contrasts the hybrid against the price-only baseline under the same execution model, complemented by ablations, drift controls, multiple-comparisons checks and paired statistical tests. The hybrid strategy did not outperform the baseline in the tested configuration. The main contributions are a reproducible negative result showing that apparent profitability does not necessarily imply predictive ability, and a reusable methodology that subsequent work can extend with richer text sources, additional instruments and broader validation protocols.

   \vspace{\onelineskip}

   \noindent
   \textbf{Keywords}: algorithmic trading. forex. event features. GDELT. backtesting. reproducibility. financial machine learning.
 \end{otherlanguage*}
\end{resumo}
